% GENERATED FILE. Edit canonical lessons, then run npm run generate:books.
% canonical-source-hash: fnv1a64:4c1c1e274c920a15
% canonical-lessons: ZH-C264-jujue, ZH-C264-kexi, ZH-C264-baoqian, ZH-C264-kongpa, ZH-C264-qishi

\chapter{Sorry, I Can't}
\label{ch:zh-sorry-i-can-t}
\hlchaptermodality{264}

\begin{chapteropening}
\textbf{By the end of this chapter:} \emph{I can turn something down politely, say it is a shame, say sorry, say I am afraid not, and say what is actually so.}

Complete the last lesson of chapter 264: I can turn something down politely, say it is a shame, say sorry, say I am afraid not, and say what is actually so.
\end{chapteropening}

\section[jùjué]{jùjué — to refuse; to turn down}

\label{lesson:ZH-C264-jujue}



\begin{quote}
Before the new one: say the Mandarin for sick leave, then the Mandarin for annual leave.
\end{quote}

\subsection*{You'll want to know: jùjué}
\emph{jùjué} — \textquotedblleft{}to refuse; to turn down\textquotedblright{}.

\begin{quote}
\emph{tā jùjuéle wǒ} — He turned me down.
\end{quote}

Its characters come later; for now it is a word for the ear.

\begin{cousinweb}
\emph{jùjué} is built from \emph{jù} 'to resist' and \emph{jué} 'to cut off', both inherited from Middle Chinese, \emph{gjoX} and \emph{dzjwet} (Guangyun, in Baxter's transcription), which glosses the first 'to resist' and the second 'to break off'.
\end{cousinweb}

\begin{grammarlens}[title={What you've built}]
One more word for saying no and asking.
\end{grammarlens}

\subsection*{Your turn}
\begin{itemize}
\raggedright
  \item \emph{Say it:} \emph{jùjué}
  \item \emph{Say it:} \emph{jùjué}, once more
  \item \emph{Say it:} \emph{tā jùjuéle wǒ} — He turned me down.
  \item \emph{Recall:} say the Mandarin for sick leave, then the Mandarin for annual leave, then say \emph{jùjué} again
\end{itemize}

\begin{tcolorbox}[breakable,colback=teal!4,colframe=teal!35!black,title={Before you move on}]
What does jùjué mean? (\textquotedblleft{}To refuse; to turn down\textquotedblright{}.) Say it once more.
\end{tcolorbox}
\section[kěxī]{kěxī — a pity; it's a shame}

\label{lesson:ZH-C264-kexi}



\begin{quote}
Before the new one: say the Mandarin for annual leave, then the Mandarin for to turn down.
\end{quote}

\subsection*{You'll want to know: kěxī}
\emph{kěxī} — \textquotedblleft{}a pity; it's a shame\textquotedblright{}.

\begin{quote}
\emph{nǐ bù néng lái tài kěxī le} — It's a real shame you can't come.
\end{quote}

Its characters come later; for now it is a word for the ear.

\begin{cousinweb}
\emph{kěxī} is built from \emph{kě} and \emph{xī}, both inherited from Middle Chinese, \emph{khaX} and \emph{sjek} (Guangyun, in Baxter's transcription), which glosses the first 'to permit' and the second 'to grudge' and, after an older dictionary, 'to grieve'.
\end{cousinweb}

\begin{grammarlens}[title={What you've built}]
One more word for saying no and asking.
\end{grammarlens}

\subsection*{Your turn}
\begin{itemize}
\raggedright
  \item \emph{Say it:} \emph{kěxī}
  \item \emph{Say it:} \emph{kěxī}, once more
  \item \emph{Say it:} \emph{nǐ bù néng lái tài kěxī le} — It's a real shame you can't come.
  \item \emph{Recall:} say the Mandarin for annual leave, then the Mandarin for to turn down, then say \emph{kěxī} again
\end{itemize}

\begin{tcolorbox}[breakable,colback=teal!4,colframe=teal!35!black,title={Before you move on}]
What does kěxī mean? (\textquotedblleft{}A pity; it's a shame\textquotedblright{}.) Say it once more.
\end{tcolorbox}
\section[bàoqiàn]{bàoqiàn — sorry; regretful}

\label{lesson:ZH-C264-baoqian}



\begin{quote}
Before the new one: say the Mandarin for to turn down, then the Mandarin for a shame.
\end{quote}

\subsection*{You'll want to know: bàoqiàn}
\emph{bàoqiàn} — \textquotedblleft{}sorry; regretful\textquotedblright{}.

\begin{quote}
\emph{hěn bàoqiàn wǒ míngtiān méiyǒu shíjiān} — I'm very sorry, I don't have time tomorrow.
\end{quote}

Its characters come later; for now it is a word for the ear.

\begin{cousinweb}
\emph{bàoqiàn} is built from \emph{bào} 'to hold' and \emph{qiàn}, both inherited from Middle Chinese, \emph{bawX} and \emph{kheamH} (Guangyun, in Baxter's transcription), which glosses the first 'to hold'. For the second it gives a sunken mouth, and in other readings 'not eating your fill'; the sense of being sorry is not given there, so how it came is not settled.
\end{cousinweb}

\begin{grammarlens}[title={What you've built}]
One more word for saying no and asking.
\end{grammarlens}

\subsection*{Your turn}
\begin{itemize}
\raggedright
  \item \emph{Say it:} \emph{bàoqiàn}
  \item \emph{Say it:} \emph{bàoqiàn}, once more
  \item \emph{Say it:} \emph{hěn bàoqiàn wǒ míngtiān méiyǒu shíjiān} — I'm very sorry, I don't have time tomorrow.
  \item \emph{Recall:} say the Mandarin for to turn down, then the Mandarin for a shame, then say \emph{bàoqiàn} again
\end{itemize}

\begin{tcolorbox}[breakable,colback=teal!4,colframe=teal!35!black,title={Before you move on}]
What does bàoqiàn mean? (\textquotedblleft{}Sorry; regretful\textquotedblright{}.) Say it once more.
\end{tcolorbox}
\section[kǒngpà]{kǒngpà — I'm afraid that; probably}

\label{lesson:ZH-C264-kongpa}



\begin{quote}
Before the new one: say the Mandarin for a shame, then the Mandarin for sorry to say.
\end{quote}

\subsection*{You'll want to know: kǒngpà}
\emph{kǒngpà} — \textquotedblleft{}I'm afraid that; probably\textquotedblright{}.

\begin{quote}
\emph{kǒngpà tā bú huì lái le} — I'm afraid he won't come.
\end{quote}

Its characters come later; for now it is a word for the ear.

\begin{cousinweb}
\emph{kǒngpà} is built from \emph{kǒng} and \emph{pà}, both inherited from Middle Chinese, \emph{khjowngX} and \emph{phaeH} (Guangyun, in Baxter's transcription), which glosses both 'to fear'.
\end{cousinweb}

\begin{grammarlens}[title={What you've built}]
One more word for saying no and asking.
\end{grammarlens}

\subsection*{Your turn}
\begin{itemize}
\raggedright
  \item \emph{Say it:} \emph{kǒngpà}
  \item \emph{Say it:} \emph{kǒngpà}, once more
  \item \emph{Say it:} \emph{kǒngpà tā bú huì lái le} — I'm afraid he won't come.
  \item \emph{Recall:} say the Mandarin for a shame, then the Mandarin for sorry to say, then say \emph{kǒngpà} again
\end{itemize}

\begin{tcolorbox}[breakable,colback=teal!4,colframe=teal!35!black,title={Before you move on}]
What does kǒngpà mean? (\textquotedblleft{}I'm afraid that; probably\textquotedblright{}.) Say it once more.
\end{tcolorbox}
\section[qíshí]{qíshí — actually; as a matter of fact}

\label{lesson:ZH-C264-qishi}



\begin{quote}
Before the new one: say the Mandarin for sorry to say, then the Mandarin for afraid that.
\end{quote}

\subsection*{You'll want to know: qíshí}
\emph{qíshí} — \textquotedblleft{}actually; as a matter of fact\textquotedblright{}.

\begin{quote}
\emph{qíshí wǒ hěn máng} — Actually, I'm really busy.
\end{quote}

Its characters come later; for now it is a word for the ear.

\begin{cousinweb}
\emph{qíshí} is built from \emph{qí} 'its' and \emph{shí} 'real', both inherited from Middle Chinese, \emph{gi} and \emph{zyit} (Guangyun, in Baxter's transcription), which calls the first a particle and glosses the second 'full' and 'true'.
\end{cousinweb}

\begin{grammarlens}[title={What you've built}]
One more word for saying no and asking.
\end{grammarlens}

\subsection*{Your turn}
\begin{itemize}
\raggedright
  \item \emph{Say it:} \emph{qíshí}
  \item \emph{Say it:} \emph{qíshí}, once more
  \item \emph{Say it:} \emph{qíshí wǒ hěn máng} — Actually, I'm really busy.
  \item \emph{Recall:} say the Mandarin for sorry to say, then the Mandarin for afraid that, then say \emph{qíshí} again
\end{itemize}

\begin{tcolorbox}[breakable,colback=teal!4,colframe=teal!35!black,title={Before you move on}]
What does qíshí mean? (\textquotedblleft{}Actually; as a matter of fact\textquotedblright{}.) Say it once more.
\end{tcolorbox}
